\begin{lemma}[Lemma~3.4, with the quasi-triangle constant carried]
Let $\mathbb G$ be a sub-elliptic setting (\noderef{CarnotSetting}) with group law $*$,
homogeneous norm $\|\cdot\|$ and $d(x,y):=\|y^{-1}*x\|$, and let
$\Om\subset\mathbb R^{n}$ be a bounded open set. Let $u,v:\mathbb R^{n}\to\mathbb R$ be
continuous on $\bar\Om$, let $\om_u,\om_v:\mathbb R\to\mathbb R$ be non-decreasing with
\[
 |u(a)-u(b)|\le\om_u\big(d(a,b)\big),\qquad |v(a)-v(b)|\le\om_v\big(d(a,b)\big)
 \qquad(a,b\in\bar\Om),
\]
let $\tau>0$, assume $u-v\le-\tau$ on $\del\Om$, let $x_M\in\bar\Om$ be a point at which
$u-v$ attains its maximum over $\bar\Om$ and assume $\tau\le M_0:=u(x_M)-v(x_M)$. Let
$\rho_0>0$ satisfy $\om_u(\rho_0)+\om_v(\rho_0)\le\tau/4$. Write $\Om_\delta$ for the inner
domain (\noderef{TranslatedDomain}), $K_\rho$ for the compact core
(\noderef{BoundaryCore}), $M_\delta(h)$ for the translated maximum value
(\noderef{MaxTranslate}),
\[
 E:=\big\{y\in\Om:\ u-v\ \text{attains its maximum over}\ \bar\Om\ \text{at}\ y\big\}
 \ \ \big(=\argmax_\Om(u-v)\big),
\]
and $u_h:=u(\cdot*h)$. Then there is $c'\ge1$ such that, setting
\[
 \rho_1:=\frac{\rho_0}{2c'},\qquad \rho_2:=\frac{\rho_0}{4c'^{2}},\qquad
 \delta_1:=\frac{\rho_0}{4c'^{2}},
\]
the following hold.
\begin{enumerate}
\item[(a)] $E\subseteq K_{\rho_0}$; and for every $\delta$ with $0<\delta\le\delta_1$,
  $M_\delta(0)=M_0$ and
  $\{x\in\Om_\delta:\ u_0-v\ \text{is maximal over}\ \Om_\delta\ \text{at}\ x\}=E$.
\item[(a$'$)] For every $\delta$ with $0<\delta\le\delta_1$, every $x\in K_{\rho_1}$ and
  every $g$ with $\|g\|<\delta$, one has $x*g\in K_{\rho_2}$.
\item[(b)] For every $\delta$ with $0<\delta\le\delta_1$ and every $h$ with $\|h\|<\delta$,
  $u_h-v\le-\tau/2$ on $\Om_\delta\mns K_{\rho_1}$.
\item[(c)] For every $\delta$ with $0<\delta\le\delta_1$ and every $h$ with $\|h\|<\delta$:
  $M_\delta(h)\ge M_0-\tau/4$; the supremum defining $M_\delta(h)$ is attained, i.e. there is
  $x\in\Om_\delta$ at which $u_h-v$ is maximal over $\Om_\delta$ and
  $M_\delta(h)=u(x*h)-v(x)$; and every such $x$ satisfies $x\in K_{\rho_1}$ and
  $x*h\in K_{\rho_2}$.
\end{enumerate}
\end{lemma}
\begin{proof}
This is Lemma~3.4 together with the sentence of (3.11) that fixes $\delta_1$.

\medskip\noindent\textbf{Why the constants are explicit, and why that is not a deviation.}
The paper sets $\delta_1:=\rho_0/4$ in (3.11) and then argues in (b) "$\dist(x,\del\Om)<\rho_0/2$,
so $\dist(x*h,\del\Om)<\rho_0/2+\delta<\rho_0$" and in (c) "$\dist(y*h^{-1},\del\Om)\ge\rho_0-\delta>\delta$".
Both are steps with triangle constant $1$. But the homogeneous norm (2.15) is a
\emph{quasi}-norm: (2.17) itself is stated as $\|x*y\|\le c_{\mathbb G}(\|x\|+\|y\|)$
with $c_{\mathbb G}=c(\mathbb G,\nu)$ for $\|x\|,\|y\|<\nu$, and that is exactly what
\noderef{CarnotSetting} provides (pm-0013 is the governing constraint). So the radii have to
carry $c_{\mathbb G}$. We therefore produce $c'\ge1$ dominating $c_{\mathbb G}$ on a scale
large enough for every application below, and replace the paper's $\rho_0/2$, $\rho_0/4$ by
$\rho_1=\rho_0/(2c')$ and $\rho_2=\delta_1=\rho_0/(4c'^{2})$. Nothing in any \emph{statement}
changes: $\rho_0$ still satisfies (3.11) verbatim, and (a), (b), (c) are the paper's
assertions with the paper's $\rho_0/2$, $\rho_0/4$ replaced by the two constants it would
have produced had (2.17) been applied with its own constant. This is the one place in
the formalization where the paper's constants are made explicit; the paper is read
precisely, not corrected. Part (a) uses no triangle inequality at all, and we keep it that
way so that the reader can see where the constants enter and where they do not.

\medskip\noindent\textbf{Step 0: preliminaries.} Since $x_M\in\bar\Om$, $\bar\Om\ne\emptyset$.
Taking $a=b=x_M$ in the modulus hypothesis for $u$ and using
$d(x_M,x_M)=\|(-x_M)*x_M\|=\|0\|=0$ (\noderef{CarnotSetting}) gives $0\le\om_u(0)$, hence
$\om_u(s)\ge0$ for every $s\ge0$ by monotonicity, and likewise for $\om_v$. Combined with
$\om_u(\rho_0)+\om_v(\rho_0)\le\tau/4$ this gives
\[
 0\le\om_u(\rho_0)\le\tau/4,\qquad 0\le\om_v(\rho_0)\le\tau/4. \tag{1}
\]
Also $M_0\ge\tau>0$, and $u-v\le M_0$ on $\bar\Om$ by the maximality at $x_M$. Since $\Om$
is open, $\del\Om=\bar\Om\mns\Om$.

\emph{$E\ne\emptyset$, indeed $x_M\in E$.} If $x_M\in\del\Om$ then $M_0=u(x_M)-v(x_M)\le-\tau<0$,
contradicting $M_0\ge\tau>0$. So $x_M\in\bar\Om\mns\del\Om=\Om$, and $u-v$ is maximal over
$\bar\Om$ at $x_M$ by hypothesis; thus $x_M\in E$. Note also
\[
 E=\{y\in\Om:\ u(y)-v(y)=M_0\}, \tag{2}
\]
because $u-v\le M_0$ on $\bar\Om$, so "maximal over $\bar\Om$ at $y$" is the same as
"$(u-v)(y)=M_0$" for $y\in\Om\subseteq\bar\Om$.

\medskip\noindent\textbf{Step 1: the scale $\nu$ and the constant $c'$.} The group law $*$ is $C^{\infty}$ and
the homogeneous norm $\|\cdot\|$ is continuous (\noderef{CarnotSetting}), and
$x\mapsto-x$ is continuous, so $(z,w)\mapsto d(z,w)=\|(-w)*z\|$ is a composition of
continuous maps and is continuous on $\mathbb R^{n}\times\mathbb R^{n}$. The set $\bar\Om$ is compact ($\Om$ is bounded) and
non-empty, so $\bar\Om\times\bar\Om$ is compact and non-empty and
\[
 \Delta:=\max\{d(z,w):\ z,w\in\bar\Om\}
\]
is a finite real number, with $\Delta\ge0$. Put $\nu:=\Delta+\rho_0+1>0$. By the
quasi-triangle inequality (2.17) of \noderef{CarnotSetting}, applied at this $\nu$,
there is $c_{\mathbb G}>0$ with
\[
 \|a*b\|\le c_{\mathbb G}\big(\|a\|+\|b\|\big)\qquad\text{whenever }\|a\|<\nu,\ \|b\|<\nu .
\]
Set $c':=\max\{c_{\mathbb G},1\}$, so $c'\ge1$, $c'\ge c_{\mathbb G}$ and, since
$\|a\|+\|b\|\ge0$,
\[
 \|a*b\|\le c'\big(\|a\|+\|b\|\big)\qquad\text{whenever }\|a\|<\nu,\ \|b\|<\nu . \tag{3}
\]
With this $c'$ define $\rho_1,\rho_2,\delta_1$ as in the statement. Then
\[
 0<\delta_1=\rho_2=\frac{\rho_0}{4c'^{2}}\le\frac{\rho_0}{4},\qquad
 \rho_1=\frac{\rho_0}{2c'}\le\frac{\rho_0}{2},\qquad
 \delta_1=\frac{\rho_1}{2c'}\le\frac{\rho_1}{2}<\rho_1<\rho_0 . \tag{4}
\]
In particular $\delta\le\delta_1<\rho_1<\rho_0<\nu$ for every admissible $\delta$, and
$\|h\|<\delta<\nu$ for every admissible $h$.

We also record once and for all the following, used repeatedly below. \emph{For every
$\delta>0$, every $x\in\Om_\delta$ and every $g$ with $\|g\|<\delta$ one has
$x*g\in\Om\subseteq\bar\Om$.} Indeed, by associativity and $0*g=g$
(\noderef{CarnotSetting}),
$d(x*g,x)=\|(-x)*(x*g)\|=\|((-x)*x)*g\|=\|0*g\|=\|g\|<\delta$ --- this is the exact
identity (2.16), kept inline by pm-0012 --- so $x*g$ lies in the $d$-ball of radius
$\delta$ about $x$, which is contained in $\Om$ by item (6) of
\noderef{InnerDomainProperties}. \hfill(4$'$)

\medskip\noindent\textbf{Step 2: the boundary estimate.} \emph{If $z\in\bar\Om$ and
$y\in\del\Om$ satisfy $d(z,y)\le\rho_0$, then $(u-v)(z)\le-3\tau/4$.} Indeed $y\in\del\Om\subseteq\bar\Om$,
so the moduli apply to the pair $(z,y)$, and using $u-v\le-\tau$ at $y$, monotonicity of
$\om_u,\om_v$ and (1),
\[
 (u-v)(z)=(u-v)(y)+\big(u(z)-u(y)\big)-\big(v(z)-v(y)\big)
 \le-\tau+\om_u\big(d(z,y)\big)+\om_v\big(d(z,y)\big)\le-\tau+\frac\tau4=-\frac{3\tau}4 .
 \tag{5}
\]

\medskip\noindent\textbf{Step 3: part (a).} \emph{$E\subseteq K_{\rho_0}$.} Let $y\in E$ and
$y'\in\del\Om$. We must show $\rho_0\le d(y,y')$ (membership $y\in\Om$ is part of the
definition of $E$). Suppose $d(y,y')<\rho_0$. Then $y\in\Om\subseteq\bar\Om$ and
$d(y,y')\le\rho_0$, so (5) gives $(u-v)(y)\le-3\tau/4<0<\tau\le M_0$, contradicting
$(u-v)(y)=M_0$ from (2). Hence $\rho_0\le d(y,y')$, and $y\in K_{\rho_0}$ by
\noderef{BoundaryCore}. No triangle inequality was used.

\emph{$M_\delta(0)=M_0$ and $A_\delta(0)=E$, for $0<\delta\le\delta_1$.} Since
$\delta\le\delta_1<\rho_0$ by (4), item (5) of \noderef{InnerDomainProperties} gives
$K_{\rho_0}\subseteq\Om_\delta$, so by the previous paragraph
\[
 E\subseteq K_{\rho_0}\subseteq\Om_\delta\subseteq\Om\subseteq\bar\Om , \tag{6}
\]
the third inclusion being item (1) of \noderef{InnerDomainProperties}. Now $u_0=u$, because
$x*0=x$ (\noderef{CarnotSetting}). For $x\in\Om_\delta$ we have $x\in\bar\Om$ by (6), hence
$u_0(x)-v(x)=(u-v)(x)\le M_0$; and $x_M\in E\subseteq\Om_\delta$ with
$u_0(x_M)-v(x_M)=M_0$. So $u_0-v$ is bounded above on the non-empty set $\Om_\delta$ by
$M_0$ and attains the value $M_0$ there; therefore its supremum over $\Om_\delta$, which is
$M_\delta(0)$ by \noderef{MaxTranslate}, equals $M_0$, and the points of $\Om_\delta$ at
which $u_0-v$ is maximal over $\Om_\delta$ are exactly the points $x\in\Om_\delta$ with
$(u-v)(x)=M_0$. By (2) and (6) that set is $E$: if $x\in\Om_\delta$ and $(u-v)(x)=M_0$ then
$x\in\Om$ by (6), so $x\in E$ by (2); conversely every $y\in E$ lies in $\Om_\delta$ by (6)
and has $(u-v)(y)=M_0$. This is part (a).

\medskip\noindent\textbf{Step 4: part (b).} Let $0<\delta\le\delta_1$, $\|h\|<\delta$, and
let $x\in\Om_\delta\mns K_{\rho_1}$. By \noderef{BoundaryCore} and $x\in\Om$, failing
$x\in K_{\rho_1}$ means there is $y\in\del\Om$ with
\[
 d(x,y)<\rho_1 . \tag{7}
\]
By (4$'$),
\[
 x*h\in\Om\subseteq\bar\Om . \tag{8}
\]
Now estimate $d(x*h,y)$. Using associativity,
$(-y)*(x*h)=\big((-y)*x\big)*h$, and $\|(-y)*x\|=d(x,y)<\rho_1<\nu$ by (7) and (4), while
$\|h\|<\delta<\nu$; so (3) applies and, with (7), (4) and $c'\ge1$,
\[
 d(x*h,y)=\big\|\big((-y)*x\big)*h\big\|\le c'\big(d(x,y)+\|h\|\big)
 <c'\Big(\frac{\rho_0}{2c'}+\frac{\rho_0}{4c'^{2}}\Big)
 =\frac{\rho_0}2+\frac{\rho_0}{4c'}\le\frac{\rho_0}2+\frac{\rho_0}4=\frac{3\rho_0}4<\rho_0 .
 \tag{9}
\]
This is the step where the paper writes $\dist(x*h,\del\Om)<\rho_0/2+\delta<\rho_0$; the
constants $\rho_1$ and $\delta_1$ are chosen precisely so that the conclusion
$d(x*h,y)<\rho_0$ survives the factor $c'$.

By (8), (9) and Step 2 applied to $z:=x*h$,
\[
 (u-v)(x*h)\le-\frac{3\tau}4 . \tag{10}
\]
Finally, $x\in\bar\Om$ and $x*h\in\bar\Om$ by (6) and (8), so the modulus of $v$ applies to
the pair $(x*h,x)$ and, with $d(x*h,x)=\|h\|<\delta\le\delta_1\le\rho_0$, monotonicity and
(1),
\[
 \big|v(x*h)-v(x)\big|\le\om_v\big(\|h\|\big)\le\om_v(\rho_0)\le\frac\tau4 . \tag{11}
\]
Therefore, by (10) and (11),
\[
 (u_h-v)(x)=u(x*h)-v(x)=\big(u(x*h)-v(x*h)\big)+\big(v(x*h)-v(x)\big)
 \le-\frac{3\tau}4+\frac\tau4=-\frac\tau2 ,
\]
which is part (b).

\medskip\noindent\textbf{Step 5: part (c), the lower bound for $M_\delta(h)$.} Let
$0<\delta\le\delta_1$ and $\|h\|<\delta$, and fix $y\in E$ (possible, $x_M\in E$ by Step 0).
We claim $y*(-h)\in\Om_\delta$.

First, $y\in K_{\rho_0}$ by Step 3, so by item (7) of \noderef{InnerDomainProperties} the
$d$-ball of radius $\rho_0$ about $y$ lies in $\Om$; and
$d\big(y*(-h),y\big)=\|(-y)*\big(y*(-h)\big)\|=\|\big((-y)*y\big)*(-h)\|=\|-h\|=\|h\|<\delta<\rho_0$,
using associativity, $0*z=z$ and $\|-z\|=\|z\|$. Hence
\[
 y*(-h)\in\Om\subseteq\bar\Om . \tag{12}
\]
Second, let $y'\in\del\Om$. Since $h*(-h)=0$ and $z*0=z$ we have
$\big(y*(-h)\big)*h=y*\big((-h)*h\big)=y*0=y$, so by associativity
\[
 (-y')*y=(-y')*\Big(\big(y*(-h)\big)*h\Big)=\Big((-y')*\big(y*(-h)\big)\Big)*h .
\]
Here $\|(-y')*\big(y*(-h)\big)\|=d\big(y*(-h),y'\big)\le\Delta<\nu$ because
$y*(-h)\in\bar\Om$ by (12) and $y'\in\del\Om\subseteq\bar\Om$, and $\|h\|<\delta<\nu$; so
(3) applies, and with $\rho_0\le d(y,y')$ from Step 3,
\[
 \rho_0\le d(y,y')\le c'\Big(d\big(y*(-h),y'\big)+\|h\|\Big),
\]
whence, using $\|h\|<\delta\le\delta_1=\rho_0/(4c'^{2})$ and $c'\ge1$,
\[
 d\big(y*(-h),y'\big)\ \ge\ \frac{\rho_0}{c'}-\|h\|\ >\ \frac{\rho_0}{c'}-\frac{\rho_0}{4c'^{2}}
 \ \ge\ \frac{\rho_0}{c'}-\frac{\rho_0}{4c'}\ =\ \frac{3\rho_0}{4c'}\ \ge\ \frac{3\rho_0}{4c'^{2}}
 \ >\ \frac{\rho_0}{4c'^{2}}=\delta_1\ \ge\ \delta .
\]
So $\delta<d\big(y*(-h),y'\big)$ for every $y'\in\del\Om$, and with (12) this gives
$y*(-h)\in\Om_\delta$ by \noderef{TranslatedDomain}. This is the paper's
"$\dist(y*h^{-1},\del\Om)\ge\rho_0-\delta>\delta$, so $y*h^{-1}\in\Om_\delta$", with the
constant carried.

Next, $\big(y*(-h)\big)*h=y$ as computed above, so
\[
 (u_h-v)\big(y*(-h)\big)=u(y)-v\big(y*(-h)\big).
\]
By (12), $y\in\bar\Om$ and $y*(-h)\in\bar\Om$, and
$d\big(y*(-h),y\big)=\|h\|\le\rho_0$, so the modulus of $v$ and (1) give
$v\big(y*(-h)\big)\le v(y)+\om_v(\|h\|)\le v(y)+\tau/4$. Hence
\[
 (u_h-v)\big(y*(-h)\big)\ \ge\ u(y)-v(y)-\frac\tau4\ =\ M_0-\frac\tau4\ \ge\ \tau-\frac\tau4
 \ =\ \frac{3\tau}4\ >\ 0 , \tag{13}
\]
using (2) and $M_0\ge\tau$.

\medskip\noindent\textbf{Step 6: part (c), attainment.} Keep $\delta$ and $h$ as in Step 5.
By (4), $\delta\le\delta_1<\rho_1$, so item (5) of \noderef{InnerDomainProperties} gives
\[
 K_{\rho_1}\subseteq\Om_\delta , \tag{14}
\]
and item (4) of \noderef{InnerDomainProperties} gives that $K_{\rho_1}$ is compact (note
$\rho_1>0$ by (4)).

\emph{$u_h-v$ is continuous on $K_{\rho_1}$.} The map $z\mapsto z*h$ is continuous
(\noderef{CarnotSetting}) and, by (14) and (8), sends $K_{\rho_1}$ into $\Om\subseteq\bar\Om$;
$u$ is continuous on $\bar\Om$; so $u_h=u(\cdot*h)$ is continuous on $K_{\rho_1}$. Also
$K_{\rho_1}\subseteq\Om\subseteq\bar\Om$ by item (2) of \noderef{InnerDomainProperties}, so
$v$ is continuous on $K_{\rho_1}$. Hence so is $u_h-v$.

\emph{$K_{\rho_1}\ne\emptyset$.} The point $y*(-h)$ of Step 5 lies in $\Om_\delta$ and, by
(13), $(u_h-v)\big(y*(-h)\big)\ge3\tau/4>-\tau/2$; so by part (b) it cannot lie in
$\Om_\delta\mns K_{\rho_1}$, i.e. $y*(-h)\in K_{\rho_1}$.

Therefore $u_h-v$ attains its maximum over the non-empty compact set $K_{\rho_1}$, say at
$x_0\in K_{\rho_1}\subseteq\Om_\delta$. We claim $u_h-v$ is maximal over all of $\Om_\delta$
at $x_0$. Let $x\in\Om_\delta$. If $x\in K_{\rho_1}$ then $(u_h-v)(x)\le(u_h-v)(x_0)$ by the
choice of $x_0$. If $x\notin K_{\rho_1}$ then part (b) gives $(u_h-v)(x)\le-\tau/2$, while
$(u_h-v)(x_0)\ge(u_h-v)\big(y*(-h)\big)\ge3\tau/4>-\tau/2$ by the choice of $x_0$,
$y*(-h)\in K_{\rho_1}$ and (13); so again $(u_h-v)(x)\le(u_h-v)(x_0)$.

Consequently $u_h-v$ is bounded above on the non-empty set $\Om_\delta$ by
$(u_h-v)(x_0)$, and attains that value there, so by \noderef{MaxTranslate}
\[
 M_\delta(h)=u(x_0*h)-v(x_0) , \tag{15}
\]
which is the asserted attainment, and by (13) and the maximality of $x_0$,
\[
 M_\delta(h)=(u_h-v)(x_0)\ \ge\ (u_h-v)\big(y*(-h)\big)\ \ge\ M_0-\frac\tau4 ,
\]
which is the asserted lower bound.

\medskip\noindent\textbf{Step 7: part (a$'$).} Let $0<\delta\le\delta_1$, let
$x\in K_{\rho_1}$ and let $\|g\|<\delta$; write $h:=g$ for uniformity with the computation
below. By (14), $x\in K_{\rho_1}\subseteq\Om_\delta$, so $x*h\in\Om$ by (8).

Let $y'\in\del\Om$. Since $h*(-h)=0$ and
$z*0=z$, $(x*h)*(-h)=x$, so by associativity
\[
 (-y')*x=(-y')*\big((x*h)*(-h)\big)=\Big((-y')*(x*h)\Big)*(-h).
\]
Here $\|(-y')*(x*h)\|=d(x*h,y')\le\Delta<\nu$ because $x*h\in\bar\Om$ by (8) and
$y'\in\bar\Om$, and $\|-h\|=\|h\|<\delta<\nu$; so (3) applies, and with
$\rho_1\le d(x,y')$ (from $x\in K_{\rho_1}$, \noderef{BoundaryCore}),
\[
 \rho_1\le d(x,y')\le c'\Big(d(x*h,y')+\|h\|\Big),
\]
whence, using $\|h\|<\delta\le\delta_1=\rho_0/(4c'^{2})$ and $\rho_1=\rho_0/(2c')$,
\[
 d(x*h,y')\ \ge\ \frac{\rho_1}{c'}-\|h\|\ >\ \frac{\rho_0}{2c'^{2}}-\frac{\rho_0}{4c'^{2}}
 \ =\ \frac{\rho_0}{4c'^{2}}\ =\ \rho_2 .
\]
In particular $\rho_2\le d(x*h,y')$ for every $y'\in\del\Om$, so $x*h\in K_{\rho_2}$ by
\noderef{BoundaryCore}. This is part (a$'$); note that it used only $x\in K_{\rho_1}$ and
$\|h\|<\delta$, not any maximality of $x$, which is why it is stated separately. It is the
paper's "$\A_\delta(h)*h\subseteq K_{\rho_0/4}$", with the constant carried.

\medskip\noindent\textbf{Step 8: part (c), $A_\delta(h)\subseteq K_{\rho_1}$ and
$A_\delta(h)*h\subseteq K_{\rho_2}$.} Let $0<\delta\le\delta_1$, $\|h\|<\delta$, and let
$x\in\Om_\delta$ be any point at which $u_h-v$ is maximal over $\Om_\delta$. Then
$(u_h-v)(x)=M_\delta(h)\ge M_0-\tau/4\ge3\tau/4>-\tau/2$ by Step 6 and (13); so by part
(b), $x\in K_{\rho_1}$. Part (a$'$) applied with $g:=h$ then gives $x*h\in K_{\rho_2}$,
which completes part (c).
\end{proof}
